\documentclass[11pt]{article}
\usepackage[margin=1in]{geometry}
\usepackage{amsmath,amssymb,amsthm,booktabs,microtype,hyperref}
\hypersetup{hidelinks,pdfauthor={T},pdftitle={Transferable Exposure Escrow Under Ambiguous Agent Effects}}
\newtheorem{theorem}{Theorem}
\newtheorem{proposition}{Proposition}
\title{Transferable Exposure Escrow Under Ambiguous Agent Effects}
\author{T}
\date{}
\begin{document}
\maketitle
\begin{abstract}
When a disconnected agent receives no receipt, its effect or rights
transfer may already have succeeded. A single conservation invariant
keeps both operations backed until their outcomes are durable.
After rollback, the largest compatible old exposure determines how
much quota can be issued again. A matching certificate sharpens the
reserve when unresolved effects compete for two exclusive rights.
We check the invariant by finite state exploration and exercise
lost acknowledgments in two local SQLite stores. In the database
control, a timeout refund exceeds the exposure cap while the quota
ledger still balances. Applying the rule requires durable custody
and enforcement at the actual effect sink.
\end{abstract}

\section*{Summary and scope}
\paragraph{Question.} How can an agent enforce a physical exposure
cap while effects, transfers and recovery acknowledgments are delayed?
\paragraph{Finding.} Keep each unresolved effect or transfer backed
until its durable outcome is known. After closing the old generation,
reserve the largest old exposure consistent with the evidence and
reissue only the difference from the cap.
\paragraph{Evidence and scope.} The invariants are proved
and checked in separate finite interleaving models. A deployed
system still needs enforcement at the actual sinks and a recovery
register outside the state that can roll back.

\section{An exposure cap that survives a network failure}
A service may allow a local agent to send a bounded number of
human-facing notifications while disconnected. If an acknowledgment
does not return, the local service cannot infer that nothing happened.
The same uncertainty applies if the service moved some of its
allowance to another device. Treating either timeout as unused
capacity makes a shared limit unreliable.

Escrow allocation permits local decisions under numeric constraints
\cite{oneil}. Bounded-counter replication applies the same principle
to rights transfer between replicas \cite{balegas}. Here the spend
being limited is a physical effect at a separate sink. We distinguish
the capacity that a ledger represents from the exposure that a person
or external system has actually received.

Let $B$ be a fixed cap. Allocate free rights among shards, totaling
$B$. All charges are nonnegative. An effect request has a unique
scoped identifier and an upper charge $r$ that covers all effects
possible under that request. Its true charge $a$ is at most $r$.
A transfer ticket similarly has a unique identifier, an amount,
a source, a destination, and authenticated metadata. IDs bind all
effect or ticket parameters, including the intended recipient.
Each shard and sink has durable state that cannot silently revert
or be cloned into a second spendable copy.

\section{Two protocols with one accounting rule}
\begin{theorem}
Let $F$ be the total free rights, $R$ the total pending effect
reservations, $C$ the total settled charges, and $T$ the total
amounts in transit. Under the stated protocol,
\[
F+R+C+T=B.
\]
If $X$ is actual cumulative physical exposure, then
\[
X\le R+C\le B.
\]
\end{theorem}

The two protocols below preserve this accounting. The proof follows
after the transitions are specified.

An effect admission atomically moves $r$ from the owning shard's
free rights into a pending reservation. A sink must enforce that
reservation and atomically either commit the authorized effect
once or fence the request against future execution. A successful
commit binds the actual charge and result to the request. A fence
attests that no effect committed and none can commit later under
the same request. The shard may settle the corresponding durable
outcome at most once. Settlement removes the pending amount,
records $a$ as settled, and returns $r-a$ to free rights. A fence
settles with charge zero.

Transfer creation atomically debits the sender's free rights and
places the amount in transit. The receiving ledger may accept the
ticket once and move the amount to its free rights. Alternatively
it can durably reject the ticket, ruling out acceptance. Only a
matching rejection permits the sender to reclaim the amount, at
most once. Acceptance and rejection are mutually exclusive.
A duplicate message can return the established status. It cannot
change rights. A ledger's local view may lag the recipient's
decision. For the invariant, a transfer remains in transit exactly
until acceptance or reclaim. An accepted transfer whose acknowledgment
is missing does not make the debited rights available to the sender.

\begin{proof}
Admission moves rights from $F$ to $R$. A sink commit changes
physical exposure without moving the ledger reservation. A fence
also leaves the reservation backed until settlement. Settlement
moves $r$ out of $R$, moves $a$ into $C$, and restores $r-a$
to $F$. A ticket debit moves rights from $F$ to $T$.
Acceptance or reclaim moves the ticket amount back to $F$
at the appropriate shard. A rejection alone moves no rights.
Each transition preserves the sum and all four terms stay
nonnegative.

Every committed effect has either its full upper charge in
$R$ or its exact charge in $C$. The categories are mutually
exclusive by settlement and identifier uniqueness. A reservation
may back an uncommitted request, but it can back no more than
one logical effect. Adding the per-effect inequalities gives
the physical bound.
\end{proof}

The atomicity requirement is at the effect boundary. Recording a
key before a separate, fallible external effect can block a valid
retry. Recording it afterward can permit a duplicate. Transaction
commit addresses related agreement problems \cite{gray}. If an
external sink lacks an enforceable commit or fence, the shard
cannot use this protocol to release an ambiguous reservation.

\section{Why elapsed time cannot establish an unused right}
\begin{proposition}
Under an asynchronous channel with no deadline enforced at the
recipient, a sender cannot infer unused transfer rights from its
local timeout alone. The same holds for an unresolved effect
reservation at a separate sink.
\end{proposition}
\begin{proof}
Compare two executions in which the sender transmits the same
ticket and receives no response before timeout. In one execution
nothing has arrived at the recipient. In the other, the recipient
has accepted the ticket and its acknowledgment is delayed.
The sender has the same observations. Reclaiming in the second
execution permits both ledgers to spend the same rights. Even
the first execution is unsafe if the ticket can arrive later
and be accepted. For an effect, compare the corresponding worlds
of an undelivered request and a committed request with a lost
acknowledgment. Returning the pending amount is unsafe when the
previous request has committed or can still commit.
\end{proof}

A lease can rule out future execution when expiry is enforced
at the sink. Expiry alone cannot establish whether an earlier
effect committed. If neither a confirmed result nor a fence is
available, the protocol may strand rights for the length of the
partition. It preserves the cap while leaving that availability
question open.

\section{Reissuing quota after snapshot rollback}
A restored snapshot may show free rights that were spent after
it was taken. State-continuity systems such as Memoir address
the related problem of detecting rollback of protected state
\cite{memoir}. Here the snapshot can help identify requests,
but cannot authorize new spending. First close the old generation
at all effect sinks and rights recipients. Closure must fence old
tokens and transfers in the same durable order as acceptance,
so old physical exposure $X_{\rm old}$ is fixed before reissue.

Let $\mathcal X$ be the nonempty set of old exposures consistent
with authenticated terminal observations and a complete
authorization inventory. Let $U=\sup\mathcal X$, with
$0\le U\le B$. When request reservations are disjoint, an
exact receipt contributes its charge, an authenticated final
absence contributes zero, and an unresolved request contributes
its reservation ceiling. Summing gives a valid upper bound.
It equals $U$ when the unresolved charges can attain their
ceilings together. If the inventory is incomplete, include an
upper bound for the missing authority. The restored snapshot's
free balance does not supply such a bound.

\begin{proposition}
The maximum quota that can be reissued and spent independently
of the remaining uncertainty is $B-U$. Let $G$ be the total
quota reissued so far. As consistent terminal evidence lowers
$U$, atomically increase $G$ to $B-U$ and issue only the
increment. The accounting record and the issuance must survive
rollback together.
\end{proposition}
\begin{proof}
If $G\le B-U$, then old exposure is at most $U$ and new
exposure is at most $G$ by the new generation's escrow. Their
sum is at most $B$. For any grant larger than $B-U$, the
definition of supremum gives a compatible old exposure that
would exceed the cap if the grant were fully spent. This proves
maximality. Compatible evidence can only reduce $\mathcal X$,
so $B-U$ can only increase. Issuing the difference from the
durable total preserves the inequality, and replay of the same
or a weaker observation issues no additional rights.
\end{proof}

If the grant is issued only in quanta of size $\eta>0$, the
largest issuable total is $\eta\lfloor(B-U)/\eta\rfloor$.
The register uses that rounded total before calculating the
increment. Terminal evidence must describe the same closed
generation, cap and inventory. Combine valid observations by
intersection. An empty compatible set signals an evidence
failure and cannot authorize further issuance.

For a cap of three and reservations of two and one, suppose
the only verified terminal record says the one-unit request
did not commit. The two-unit request may have committed in full,
so recovery can safely reissue one unit. It can issue another
only after terminal evidence reduces the bound on the two-unit
request. If restoring the reissue register makes it forget the
first unit, minting that unit again can violate the cap even
though the old generation stays fenced.

This reissued unit funds an authorized fresh effect. It does
not establish that the two-unit request was absent and cannot
justify retrying that request under a new identifier. Reissue
preserves the aggregate budget. A uniqueness requirement still
needs terminal occupancy evidence for the particular effect.

The recovery model explored 67 reachable states and independently
enumerated compatible worlds for all 12 partial terminal reports.
The proposed mechanism stayed within the physical cap. Four
mutants exceeded it: treating unknown as unused, using absence
observed before the fence, issuing the full recoverable amount
twice, and rolling back the issuance register. Each witness
uses physical exposure four under cap three. This check models
the durability condition. It does not demonstrate an implementation
of a rollback-resistant register.

\section{Reissue when unresolved requests compete}
The ceilings of unresolved requests can substantially overstate
what could have committed together. Suppose each request uses
one exclusive grant key and one exclusive effect key. A trusted
old-generation sink could commit at most once at each key.
An authenticated final absence removes the associated request.
A verified commit contributes its actual charge to $K$ and
occupies both keys, even if its charge is zero. Remove those
occupied keys from the unresolved inventory. This step requires
durable evidence that the exclusivity rules held before closure.
Two verified commits on one exclusive key contradict that model.

Use a bipartite graph with grant keys on one side and effect
keys on the other. An unresolved request is an edge with its
charge ceiling $r_e$. Any jointly committed set forms a
matching. Let $M$ be the weight of a maximum matching.
Then $\min(B,K+M)$ bounds the old exposure. If the constraints
of the residual evidence are exactly the matching constraints
and the matching ceilings are attainable together under the
old cap, $K+M$ is the sharp old maximum.

A checker can certify the matching optimum without running the
solver. Supply a matching and nonnegative prices $\alpha_u$
and $\beta_v$ on the two key families. Check
$\alpha_u+\beta_v\ge r_{uv}$ for every unresolved edge.
Every matching uses each vertex at most once, so its total
weight is at most the sum of all prices. Equality with the
supplied matching proves optimality. Such a certificate exists
by linear-programming duality. The primal maximizes weighted
edge usage under unit capacity at each vertex. Negate the rows
of its incidence matrix on one side to obtain one positive and
one negative entry per column. Every square submatrix has
determinant zero or $\pm1$, by expanding a column with at most
one entry or, if every column has two, summing its rows.
The unit-capacity polytope therefore has integral vertices.
The dual variables are precisely the vertex prices.

For example, take a cap of 12 and one confirmed unit of
exposure. In the residual inventory, the first grant key
has edges of weights five and four to two effect keys. The
second grant key has one edge of weight four to the first
effect key. The two four-unit requests can commit together,
and grant prices $(4,3)$ and effect prices $(1,0)$ certify
the maximum of eight. The old exposure can reach nine and
three units can be reissued. Summing all three unresolved
ceilings would leave no quota to reissue after clipping to
the old cap.

The graph gives a safe upper bound even when additional
constraints make some matchings unattainable. For example,
two disjoint requests each with indivisible charge four and
an old cap of five can have old exposure at most four.
Clipping the unconstrained matching weight to five is safe
and conservative. The matching verifier agrees with independent
vertex-price search on 384 synthetic instances. It does not
establish that a real authorization inventory is complete.

\section{A complete finite interleaving check}
The verifier models two shards with total cap three, initial
free rights two and one, four scoped effect requests, and a
one-unit ticket in each direction. It explores all reachable
states and records ledger liability separately from physically
committed effects. Redeliveries that return existing status add
no state transition. The model includes delayed settlement and
allows a fenced effect to return its full reservation.
\begin{center}
\begin{tabular}{lrrr}
\toprule
Protocol & States & Transitions & Maximum physical exposure \\
\midrule
Backed until durable outcome & 2,592 & 6,648 & 3 \\
Effect timeout refunds rights & 12,192 & 37,936 & 5 \\
Transfer timeout reclaims rights & 8,436 & 24,256 & 5 \\
Recipient accepts a ticket twice & 6,240 & 19,184 & 5 \\
\bottomrule
\end{tabular}
\end{center}
The safe graph has neither a conservation error nor an unbacked
physical effect. Every mutant reaches a state with physical
exposure four or higher. One short transfer trace sends a
one-unit ticket, refunds it at timeout, and uses the sender's
restored balance to commit two units. The recipient then accepts
the delayed ticket and commits two more units. With cap three,
that yields a physical violation even though each shard observes
enough free rights locally.

\section{Retrying effects across separate durable stores}
The local experiment used one SQLite database for the escrow
ledger and another for effect outcomes. An effect is a committed
row at the sink. The ledger reserves the request ceiling in one
write transaction. At the sink, the request key and parameters
identify one terminal outcome, either a committed row or a fence.
Replaying the same request returns the recorded outcome. Reusing
the key with a different ceiling is rejected. Settlement matches
the receipt to the reservation and credits the difference between
its ceiling and charge in one write transaction.

The cap was five. A request reserved four and committed an effect
of charge three. A child process exited without returning the
receipt. Reopening both stores showed one free unit, four pending,
and three committed at the sink. A two-unit request still failed
admission. Retrying returned the existing receipt, which released
one unit after settlement. A second request then committed two
units, bringing exposure to five. Exiting after the first
settlement commit and retrying did not release the unit twice.

A separate exit before the sink commit left no effect. A durable
fence prevented a delayed retry from creating one. Exiting before
settling that fence left the reservation in place, and a retry
returned the full ceiling once. In the negative control, the
ledger treated the missing reply after a successful commit as
an absent effect. That allowed it to spend five additional units.
The ledger still balanced, while the sink held eight units
against the cap of five.

The fixture used SQLite 3.54.0 with write-ahead logging and five
forced process exits. It measures transactions on disposable local
rows. Real delivery requires binding the same effect identifier
and terminal record at the service that performs the action.
This experiment supplies neither that external transaction boundary
nor evidence about hardware power loss.

\section{Application and evidence boundary}
The protocol can govern budgets for outbound notifications,
sandboxed changes, or other effects with a meaningful upper
charge. Choose the charge unit before admitting requests and
log both potential liability and committed outcomes. A prospective
exercise should delay acknowledgments, interrupt settlement,
replay transfers, and crash the sink at the effect boundary.
Measure stranded capacity and safe throughput as well as cap
violations. A ledger rollback, false upper charge, or effect that
bypasses the enforcement point invalidates the bound. This
interleaving model and the local transaction experiment establish
different parts of the protocol. A deployed network still needs an
enforced upper charge and matching receipts from the real sink.

\paragraph{Review status.} Independent manuscript review is pending.

\begin{thebibliography}{9}
\bibitem{oneil} Patrick E. O'Neil. The Escrow Transactional Method.
ACM Transactions on Database Systems 11(4), 1986.
\url{https://ics.uci.edu/~cs223/papers/p405-o_neil.pdf}.
\bibitem{balegas} Valter Balegas et al. Extending Eventually
Consistent Cloud Databases for Enforcing Numeric Invariants.
2015. \url{https://arxiv.org/abs/1503.09052}.
\bibitem{gray} Jim Gray and Leslie Lamport. Consensus on
Transaction Commit. ACM Transactions on Database Systems 31(1),
2006. \url{https://arxiv.org/abs/cs/0408036}.
\bibitem{memoir} Bryan Parno, Jacob R. Lorch, John R. Douceur,
James Mickens, and Jonathan M. McCune.
Memoir: Practical State Continuity for Protected Modules.
\emph{IEEE Symposium on Security and Privacy}, 2011.
\url{https://www.jonmccune.net/papers/PLDMM2011.pdf}.
\end{thebibliography}
\end{document}
